\section{Rational certificates for the scalar dilation}
\label{app:scalar-certificate}
This appendix proves the strict bounds in \eqref{eq:cstar}. All finite
arithmetic checks below are reproduced by the standalone standard-library
script \texttt{scripts/verify\_scalar\_certificate.py}; no floating-point
evaluation is used by the certificate.

\subsection{An elasticity bound}
For $x=\rho^{-1}(y)>0$, define
\[
 E(y)=\frac{yp'(y)}{p(y)},\qquad
 K(y)=\frac{y^2(1-x^2)}{2x^2(1+x^2)}.
\]
Direct differentiation using $p=b'\circ\rho^{-1}$ gives
\begin{equation}
 \frac{\dd E}{\dd\log y}=E-K,\qquad E\geq1,
 \qquad 0<K<1.
 \label{eq:elasticity-ode}
\end{equation}
For the last inequality, set
$B(x)=\sqrt2x\sqrt{(1+x^2)/(1-x^2)}$, so $K=(\rho(x)/B(x))^2$.
With $u=x^2$, the derivative ratio is
\[
 \frac{B'}{\rho'}=\frac{1+2u-u^2}{(1+u)\sqrt{1-u}}>1,
\]
because the difference of squared numerator and denominator is
$u(3+3u-3u^2+u^3)>0$. Equality of the two functions at zero proves
$\rho<B$. Convexity of $p$ gives $E\geq1$, so $E$ increases strictly.
Solving the differential inequality in \eqref{eq:elasticity-ode} on a
multiplicative ray gives $E(ty)>1+t(E(y)-1)$ for $t>1$. Consequently,
for every $a>1$,
\begin{equation}
 \log\frac{p(ay)}{yp'(y)}>
 h_1(E):=\log a+(a-1)(E-1)-\log E.
 \label{eq:elasticity-first}
\end{equation}

For a stronger bound when $E>2$, concavity of the square root gives
\[
 \rho(x)\leq\sqrt2D(x),\qquad
 D(x)=\tfrac32\operatorname{artanh}x-x/2.
\]
At $x=4/5$, $D=(3/2)\log3-2/5<2497/2000$.
Indeed $\log3<1099/1000$ follows by summing the exponential Taylor
polynomial through degree six. Hence
$E(\rho(4/5))<2497\sqrt{41}/8000<2$, the latter inequality being
$2497^2\cdot41<64\cdot2000^2$. Thus $E(y)>2$ implies $x>4/5$.
On this region,
\[
 K\leq G(x):=\frac{D(x)^2(1-x^2)}{x^2(1+x^2)}.
\]
The positive power series for $\operatorname{artanh}$ gives
$D(x)\geq x(1+x^2/2)$, and
$D'(x)=(1+x^2/2)/(1-x^2)$. It follows that
\[
 \tfrac12(\log G)'=
 \frac{D'}D-\frac{x}{1-x^2}-\frac1x-\frac{x}{1+x^2}
 \leq-\frac{x}{1+x^2}<0.
\]
Therefore
\[
 K<G(4/5)<\left(\frac{2497}{2000}\right)^2\frac{225}{656}
 =\frac{56115081}{104960000}<k:=\frac{107}{200}.
\]
This bound holds along the entire multiplicative ray starting at $E>2$.
Integrating \eqref{eq:elasticity-ode} now gives
$E(ty)>k+t(E(y)-k)$ for $t>1$, and hence
\begin{equation}
 \log\frac{p(ay)}{yp'(y)}>h_k(E):=
 k\log a+(a-1)(E-k)-\log E\qquad(E>2).
 \label{eq:elasticity-second}
\end{equation}

\subsection{The strict upper certificate}
Take $a=69/50$. On $1\leq E\leq2$, $h_1$ is decreasing, and
\[
 h_1(E)\geq h_1(2)=\log(69/100)+19/50>0,
\]
because
$e^{19/50}>1+19/50+(19/50)^2/2=7261/5000>100/69$.
On $E>2$, the convex function $h_k$ has its global minimum at $E=50/19$,
where its value is
\[
 \log(19/50)+1+\frac{107}{200}\bigl(\log(69/50)-19/50\bigr).
\]
The positive series
$\log((1+z)/(1-z))=2\sum_{j\geq0}z^{2j+1}/(2j+1)$ yields, with
$r_0=31/69$ and $s_0=19/119$,
\begin{align*}
 \log(19/50)&>
 -2\left(r_0+\frac{r_0^3}3+\frac{r_0^5}{5(1-r_0^2)}\right)
 =-\frac{9064603453}{9362506500},\\
 \log(69/50)&>2(s_0+s_0^3/3)=\frac{1628072}{5055477}.
\end{align*}
Substitution leaves the strictly positive rational lower bound
\begin{equation}
 h_k(50/19)>
 \frac{255839420916449}{315546241820670000}>0.
 \label{eq:upper-certificate-margin}
\end{equation}
Thus $p((69/50)y)>yp'(y)$ for every $y>0$. Since the maximum defining
$c_\star$ is attained, the upper bound on $c_\star$ is strict.

\subsection{The strict lower certificate}
The elementary transform in \eqref{eq:scalar-elementary} has series
\begin{equation}
 Y(z)=\sum_{j\geq0}\frac{(2-2^{-j})z^{2j+1}}{2j+1},\qquad
 S_N(z)<Y(z)<S_N(z)+\frac{2z^{2N+3}}{(2N+3)(1-z^2)}.
 \label{eq:Y-rational-series}
\end{equation}
Here $S_N$ includes $j=0,\ldots,N$; the tail estimate follows by bounding
$2-2^{-j}<2$ and every tail denominator by $2N+3$.
Set
\[
 v_0=\frac{3787}{4000},\quad w_0=\frac{9797031}{10^7},\quad
 a_0=\frac{68743}{50000}.
\]
Summation through $N=1200$ and rational squaring for $A$ verify
\begin{gather*}
 \frac{2453756741730599}{10^{15}}<Y(v_0)
 <\frac{12268783708653}{5\cdot10^{12}},\qquad
 \frac{134943211257327}{4\cdot10^{13}}<Y(w_0),\\
 A(v_0)>\frac{405367307458211}{4\cdot10^{13}},\qquad
 A(w_0)<\frac{25381942601625079}{10^{15}}.
\end{gather*}
These finite enclosures imply
\begin{align}
 Y(w_0)-a_0Y(v_0)&>
 \frac{2071874360571}{250000000000000000}>0,\notag\\
 Y(v_0)A(v_0)-w_0A(w_0)&>
 \frac{2049519399569150216498389}{40000000000000000000000000000}>0.
 \label{eq:lower-certificate-margins}
\end{align}
The second inequality states $P(w_0)<Y(v_0)A(v_0)$.
Since $P$ and $Y$ are strictly increasing, it follows that
$Y(W(v_0))>Y(w_0)>a_0Y(v_0)$. This is a strictly feasible witness for
$c_\star>a_0$. The interval bounds, tail estimate, and rational squaring
are a finite proof certificate, independent of a numerical optimizer or
an assumption about stationary-point uniqueness.

\subsection{Every scalar-dilation maximizer lies in a certified interval}
The same inequalities also locate every maximizer of the distinct
standard-barrier constant $c_\star$. Keep $a_0=68743/50000$
and $k=107/200$ in \cref{eq:elasticity-first,eq:elasticity-second}.
Exact positive-series bounds for logarithms give
\begin{align*}
 h_1(2)&>2589087717927/40000000000000000>0,\\
 h_k(571/250)&>2011885403423/100000000000000000>0,\\
 h_k(773/250)&>8579961311243/500000000000000000>0.
\end{align*}
For clarity these are finite certificates: for rational $q\in(0,1)$,
the partial sum of $2\sum_{j\geq0}q^{2j+1}/(2j+1)$ through
$j=N$ bounds $\log((1+q)/(1-q))$ below, and adding
$2q^{2N+3}/[(2N+3)(1-q^2)]$ bounds it above. Taking
$q=(z-1)/(z+1)$ for $z=a_0,2,571/250,773/250$ and $N=200$
verifies the displayed inequalities using rational arithmetic.
The function $h_1$ decreases on $[1,2]$. The convex function
$h_k$ decreases until $E=1/(a_0-1)$ and increases thereafter;
this minimum lies strictly between $571/250$ and $773/250$.
Therefore a scalar ratio greater than $a_0$ is possible only for
$571/250<E<773/250$.

In the coordinate of \eqref{eq:scalar-elementary}, $E=Y(v)/v$
is strictly increasing by \eqref{eq:elasticity-ode}.
The series \eqref{eq:Y-rational-series} through $N=1200$ gives
\begin{align*}
 \frac{571}{250}\frac{4611}{5000}-Y(4611/5000)
 &>159652343530451/200000000000000000>0,\\
 Y(9701/10000)-\frac{773}{250}\frac{9701}{10000}
 &>19492543073837/250000000000000000>0.
\end{align*}
Since every global maximizer has ratio $c_\star>a_0$, these
inequalities, followed by rational squaring in
$x=v/\sqrt{2-v^2}$, prove
\begin{equation}
 \frac{4611}{5000}<v_\star<\frac{9701}{10000},\qquad
 \frac{43}{50}<x_\star<\frac{943}{1000}.
 \label{eq:cstar-maximizer-location}
\end{equation}
The verification script includes every displayed bound. Finally,
differentiating $P(W(v))=Y(v)A(v)$ shows that every stationary
point satisfies
\[
 Y(v)[Y'(v)A(v)+Y(v)A'(v)]
 =Y(W(v))Y'(v)A(W(v)),\qquad
 A'(v)=\frac{v(3-v^2)}{(1-v^2)^2\sqrt{2-v^2}}.
\]
No uniqueness of this stationary point or of the $c_\star$
maximizer is inferred from its certified localization.
